\chapter{Related Work}
\label{chap:related}

This chapter situates the thesis in four bodies of work. I cover explicit 3D scene representations and where 3DGS sits among them; latent generative modelling of 3D content, with attention to the two-stage VAE+generator recipe; learning on unordered sets and the use of space-filling curves to impose a canonical order; and contrastive plus disentangled representation learning, which informs my small-scale semantic experiments. Each section closes by stating the specific gap the thesis addresses.

%% ──────────────────────────────────────────────────────────────────────────
\section{3D Scene Representations}
\label{sec:rw_repr}

\paragraph{Neural radiance fields.}
NeRF \citep{mildenhall2020nerf} represents a scene as a continuous volumetric function encoded in the weights of a coordinate-conditioned MLP. It produces photorealistic novel views and remained the dominant approach for several years. Two limitations hold it back as a generative target: rendering is slow because each ray needs hundreds of network evaluations, and the per-scene weight representation has no explicit parameter for a network to predict on a scene level.

\paragraph{3D Gaussian Splatting.}
3DGS \citep{kerbl2023gaussian} represents a scene as a set of anisotropic Gaussian primitives rasterised by alpha compositing. The approach renders in real time, matches NeRF on novel-view synthesis, and exposes an explicit per-primitive parameterisation. Each primitive has a position $\mu \in \bR^3$, a colour $c \in \bR^3$ (the DC spherical-harmonic coefficient, ignoring higher-order view-dependent terms here), an opacity $\alpha \in (0,1)$, a per-axis scale $s \in \bR^3_+$, and a unit quaternion $q$ defining the orientation. The covariance matrix is then $\Sigma = R(q)\,\mathrm{diag}(s^2)\,R(q)^\top$.

Scaffold-GS \citep{lu2024scaffold} reduces redundancy in a dense 3DGS fit by anchoring per-primitive attributes to a sparser neural scaffold, predicting offsets relative to anchor points rather than absolute positions. This anchor-plus-offset structure is what motivates the position decoding choice in \S\ref{sec:method_anchor}. GaussianCube \citep{zhang2024gaussiancube} takes a different route: it fits a fixed Gaussian count at the object level and then transports the primitives onto a regular voxel grid, so a standard convolutional diffusion model can operate on a 3D tensor. Both of these object-level approaches inform my own per-token anchor decoding, but neither extends directly to unbounded indoor scenes.
\paragraph{Gap.} 3DGS gives an explicit representation, but turning it into a generative target at scene level requires a compact, organised latent. Neither 3DGS itself nor Scaffold-GS provides that latent; building it is the goal of Stage~1 in this thesis.

%% ──────────────────────────────────────────────────────────────────────────
\section{Latent Generative Modelling of 3D Content}
\label{sec:rw_gen}

\paragraph{The two-stage recipe.}
A widely used recipe for high-resolution generation is to first train a VAE to compress the signal into a latent \citep{kingma2014vae}, and then train a generative model in that latent space \citep{rombach2022ldm}. The generative half has shifted from DDPM \citep{ho2020ddpm} to flow matching \citep{lipman2023flowmatching, liu2023rectified}, both typically realised with a transformer backbone such as DiT \citep{peebles2023dit} or SiT \citep{ma2024sit}. The latent decouples the awkward structure of the raw signal from the easier learning problem of modelling a tame distribution of token sequences.

\paragraph{Object-level 3D generation.}
LION \citep{zeng2022lion} applies hierarchical latent diffusion to point cloud objects. Michelangelo \citep{zhao2023michelangelo} uses a cross-attention Perceiver \citep{jaegle2021perceiver, jaegle2022perceiverio} to encode a 3D shape into a fixed set of latent tokens aligned with image and text embeddings; its shape-latent VAE is the architectural ancestor of mine. GaussianCube \citep{zhang2024gaussiancube} structures the per-primitive prediction into a position-plus-offset form so a dense convolutional model can operate on it. All three are object-level, with bounded extent and small fixed primitive counts.

\paragraph{Scene-level 3DGS.}
Can3Tok \citep{gao2025can3tok} is the direct predecessor of this thesis. It uses a Perceiver-style cross-attention encoder to compress a scene's Gaussian primitives into a fixed-size latent, and a transformer-plus-MLP decoder to reconstruct the per-primitive parameters. Can3Tok was demonstrated on a few hundred scenes and showed conditional generation as a downstream application. SceneSplat \citep{li2025scenesplat} introduces the first generalisable feed-forward 3DGS encoder, pre-trained on its accompanying SceneSplat-7K dataset of 7{,}916 indoor 3DGS scenes; the dataset comes with per-Gaussian ScanNet-72 \citep{dai2017scannet} category labels. SceneSplat-7K is the training corpus throughout this thesis. SceneSplat++ \citep{ma2025scenesplatpp} extends the dataset and benchmarks language-aligned Gaussian Splatting more thoroughly.

\paragraph{Gap.} What is required to scale a scene-level 3DGS VAE from hundreds to several thousand pre-fitted scenes, and which architectural choices control whether it converges, has not been systematically studied. The seven-experiment ablation in Chapter~\ref{chap:results} is my answer.

%% ──────────────────────────────────────────────────────────────────────────
\section{Learning on Unordered Sets and Space-Filling Curves}
\label{sec:rw_sets}

\paragraph{Permutation-invariant encoders.}
PointNet \citep{qi2017pointnet} established symmetric aggregation as a way to produce permutation-invariant features over point sets. The Perceiver \citep{jaegle2021perceiver} generalises this by cross-attending a small learned query set against the input, yielding a fixed-size latent that is invariant to input permutation. Permutation invariance is cheap on the encoder side; the difficulty reappears on the decoder side, where any element-wise reconstruction loss compares predicted slot $i$ to target slot $i$ and so implicitly fixes an ordering.

\paragraph{Canonical orderings via space-filling curves.}
Rather than handle permutation in the loss, a complementary approach is to fix a canonical order on the data. Point Transformer V3 \citep{wu2024ptv3} serialises point cloud tokens along the Hilbert space-filling curve \citep{skilling2004hilbert} and shows that this ordering, with its tight locality property (consecutive indices remain spatially adjacent), gives better downstream results than the alternatives such as Z-order (Morton) curves, where high-order bit flips cause long spatial jumps. The same idea, but pushed onto the reconstruction \emph{target}, is the central mechanism behind RQ1 in this thesis.

\paragraph{Gap.} Canonical serialisation is well known on the processing side of 3D learning but, to the best of my knowledge, has not been used with a fixed-frame coordinate system to make the per-slot 3DGS regression target identifiable for a generative VAE. Doing this in concert with a per-token spatial anchor is what makes the slot-aligned loss learnable in my setup.

%% ──────────────────────────────────────────────────────────────────────────
\section{Semantic and Disentangled Representations}
\label{sec:rw_sem}

\paragraph{Contrastive learning.}
Self-supervised contrastive learning \citep{oord2018cpc, chen2020simclr} learns features by pulling together related samples (positives) and pushing apart unrelated ones (negatives) in embedding space. The InfoNCE objective \citep{oord2018cpc} is the standard tool. I adapt it in two forms in this thesis: a per-Gaussian variant that uses the per-Gaussian ScanNet-72 labels from SceneSplat-7K as the positive criterion, and a scene-level variant that uses the scene's label-distribution histogram.

\paragraph{Disentangled latent representations.}
$\beta$-VAE \citep{higgins2017betavae} attempts to disentangle independent factors of variation in the latent by reweighting the KL term in the VAE objective. Locatello et al.\ \citep{locatello2019challenging} show that unsupervised disentanglement is in general unidentifiable without inductive bias, which is why my small-scale experiments rely on structural inductive bias (a layout/geometry split with auxiliary heads and cross-reconstruction) rather than on KL reweighting alone.

\paragraph{Gap.} Per-Gaussian semantic labels for 3DGS scenes have become available with SceneSplat-7K \citep{li2025scenesplat}. So far they have been used for discriminative scene understanding rather than to add semantic structure to a generative latent. My 300-scene exploration in \S\ref{sec:results_300} is a first look at how that supervision interacts with a Stage~1 latent designed for downstream generation.